\cotefile{76}
\scanpage{33}
donné par $\sigma^{\unsure{\ast|\alpha}}=i_{B,A}^{\ast}$. \unsure{Quand|Faisant} varie $A$ dans $\mathrm{Pf}_{\unsure{n+1|n-1|m+1}}(\mathbb N)$, on trouve
\[ \sigma^{\ast}\colon \Theta_m(D)\unsure{\text{---}|\longrightarrow}\Theta_n(D) \]
avec transitivité évidente. Donc les $\Theta_n(D)$ $(n\in\mathbb N)$ forment un ens. quasi simplicial (op. simpliciales définies pour applications \uline{str$^{t}$. croissantes}) i.e. définissent un foncteur
\[ \underset{\substack{\text{ens. finis tot}^{\mathrm t}\\\text{ordonnés, avec}\\\text{appl. str. cr.}}}{\left(\frac12\mathrm{Simpl}^{\ast}\right)^{\circ}}\longrightarrow(\mathrm{Ens}) \]

[ Tout ceci marche en remplaçant $\mathbb N$ par un ens ordonné quelconque \inserted{$N$}, $\mathrm{Pf}(\mathbb N)$ par l'ens \inserted{$\mathrm{Dr}(N)$} des parties finies tot$^{t}$ ordonnées de $N$, i.e. l'ens des drapeaux de $N$, qui est aussi l'ens ordonné associé \deleted{au} \inserted{à la cat. fibrée définie par le} foncteur
\[ \left(\frac12\mathrm{Simpl}^{\ast}\right)\xrightarrow{\varphi_N}\mathrm{Ens} \]
donné par
\[ \varphi_N(\Delta_n)=\mathrm{Mon.}\,\mathrm{croiss}\,(\Delta_n,N) \]
\diagram{\begin{tikzcd}[ampersand replacement=\&]\mathrm{Dr}(N)\arrow[d, "{\text{--- appl. str. croissantes}}"]\\\frac12\mathrm{simpl}\end{tikzcd}}
% archive sheet number: 32
